\uuid{AvDG}
\titre{Densité}
\auteur{bodin ; hanani ; mzari}
\organisation{exo7}
\licence{CC BY-NC-SA 4.0}
\datecreate{2020-01-20}
\niveau{L1}
\module{Analyse}
\chapitre{Continuité, limite et étude de fonctions réelles}
\sousChapitre{Limite de fonctions}
\difficulte{3}
\theme{Limites des fonctions réelles, densité}
\exo7qcmid{lille-1.limites.024}

\contenu{
\begin{question}
\qclassification{123.03}
\qtags{motcle=densité}
Soit  $f$ une fonction définie sur $]0,1[$ par :
$f(x)=1$, si $x \in \Rr \setminus \Qq$ et $f(x)=\frac{1}{m},$ si $x= \frac{n}{m},$ où  $n, m \in \Nn^*$ tels que $ \frac{n}{m}$ soit une fraction irréductible. Quelles sont les assertions vraies ?
\begin{answers}

    \bad{$\lim_{x\to 1^-} f(x)=0$}

    \good{$f$ n'admet pas de limite en $1^-$.}

    \bad{$\lim_{x\to 1^-} f(x)=1$}

    \bad{$\lim_{x\to 1^-} f(x)=+\infty$}
\end{answers}
\begin{explanations}
L'ensemble des irrationnels est dense dans $\Rr$. Donc il existe une suite d'irrationnels $(u_n)$ qui tend vers $1^-$ et donc  $\lim_{n\to +\infty} f(u_n)=1$.  D'autre part, la suite $(\frac{n}{n+1})$  tend vers $1^-$ et $\lim_{n\to +\infty} f(\frac{n}{n+1})=\lim_{n\to +\infty} \frac{1}{n+1} = 0 $. On déduit que  $f$ n'admet pas de limite en $1^-$.
\end{explanations}
\end{question}
}
